% ================================================================
% SECTION 3 — METHODOLOGY
% Target: ~11–12 pages
% ================================================================

\section{Methodology: A Causal AI Framework for Time-Series Data}
\label{sec:methodology}

\subsection{Overall Framework and Workflow}
\label{sec:framework}

The analytical framework developed in this thesis follows an
eleven-step pipeline that moves from raw panel data to causal effect
estimates with formal robustness verification.
Each step is lag-aware and time-aware: no information from future
periods is allowed to influence estimates of past or present periods,
and all cross-validation splits respect chronological ordering.
Figure~\ref{fig:workflow} summarises the pipeline.

\begin{figure}[H]
  \centering
  \renewcommand{\arraystretch}{1.4}
  \begin{tabular}{cl}
    \toprule
    \textbf{Step} & \textbf{Description} \\
    \midrule
    1  & Problem formulation and treatment definition \\
    2  & Time-series data collection and corridor panel construction \\
    3  & Temporal exploratory data analysis (EDA) \\
    4  & Data preprocessing and outlier treatment \\
    5  & Ground-truth temporal graph construction \\
    6  & Five SCM-family causal estimators applied \\
    7  & Discovery evaluation: pre-period fit metrics \\
    8  & Robustness: in-time placebo and leave-one-out tests \\
    9  & Causal estimand identification \\
    10 & Causal effect estimation and confidence intervals \\
    11 & Refutation, sensitivity analysis, and reproducibility \\
    \bottomrule
  \end{tabular}
  \caption{Eleven-step causal AI framework. Steps 1--5 are
    preparatory; steps 6--8 address causal discovery and evaluation;
    steps 9--11 address causal effect estimation and validation.}
  \label{fig:workflow}
\end{figure}

\subsection{Data Description and Temporal EDA}
\label{sec:data}

\subsubsection{Data Source and Panel Structure}
\label{sec:datasource}

The dataset is a monthly balanced panel of six South African freight
rail corridors operated by Transnet Freight Rail, covering the period
January 2015 to December 2024 --- a total of 120 months per corridor
and 720 corridor-month observations.
The raw data were provided as an Excel workbook
(\texttt{NATCOR\_corridor\_panel\_raw.xlsx}) assembled from
Transnet operational records, supplemented by South African Reserve
Bank macroeconomic indices and South African Weather Service
precipitation records.

The six corridors in the panel are:
\begin{itemize}
  \item \textbf{NATCOR} --- the National Corridor (treated unit),
        connecting Gauteng to the ports of Durban and Richards Bay
        via KwaZulu-Natal.
  \item \textbf{CAPE} --- the Cape Corridor, linking Gauteng to
        the Port of Cape Town via the Western Cape.
  \item \textbf{IRON\_ORE} --- the Sishen--Saldanha iron ore
        export line, a dedicated bulk commodity corridor.
  \item \textbf{NE\_CORRIDOR} --- the North-East Corridor,
        linking Gauteng to the Mozambican border at Ressano Garcia.
  \item \textbf{NORTHCOR} --- the Northern Corridor, connecting
        Gauteng to the Zimbabwean border at Beitbridge.
  \item \textbf{SOUTHCOR} --- the Southern Corridor, serving the
        Eastern Cape ports of Port Elizabeth and East London.
\end{itemize}

The outcome variable is \textbf{Rail\_Vol\_mt}: monthly rail freight
volume in millions of tonnes (MT).
The treatment date is \textbf{April 2022} (the month of the KZN
floods), giving a pre-period of 87 months (2015-01 to 2022-03)
and a post-period of 33 months (2022-04 to 2024-12).

\subsubsection{Variables}
\label{sec:variables}

In addition to the outcome variable, the panel contains 13 covariates
measured at monthly frequency:

\begin{itemize}
  \item \textbf{Road\_Vol\_Est\_mt} --- estimated road freight volume
        on parallel routes (million tonnes); a proxy for modal
        competition and demand-side conditions.
  \item \textbf{Market\_Share\_Pct} --- rail's percentage share of
        total freight on the corridor.
  \item \textbf{Avg\_Delay\_Hrs} --- average locomotive delay per
        journey (hours); an operational performance indicator.
  \item \textbf{Loco\_Avail\_Rate} --- proportion of the locomotive
        fleet available for service (0--1); a supply-side constraint.
  \item \textbf{Load\_Shedding\_Hrs} --- monthly hours of
        load-shedding (electricity rationing) experienced in the
        region; a systemic confounder that affects both rail
        operations and economic activity.
  \item \textbf{Fuel\_Price\_Diesel} --- monthly diesel retail price
        in rands per litre; affects road freight costs and hence
        modal competition.
  \item \textbf{Regional\_GDP\_Idx} --- regional economic activity
        index (base 100 = January 2015); captures demand-side
        macroeconomic conditions.
  \item \textbf{Security\_Incidents} --- count of cable theft, rail
        vandalism, and related security incidents per month; a
        chronic operational challenge for South African freight rail.
  \item \textbf{Rainfall\_mm} --- monthly precipitation (mm);
        relevant both as a potential confounder for corridor
        operations and as a proxy for weather-related disruptions.
\end{itemize}

Table~\ref{tab:natcor_stats} presents descriptive statistics for the
NATCOR corridor over the full 120-month panel.
The pre-flood mean of \textbf{Rail\_Vol\_mt} was 3.94\,MT/month,
falling to 3.63\,MT/month in the post-flood period --- a raw
decline of 7.9\%.
Market share (NATCOR volume as a fraction of total six-corridor
throughput) fell from approximately 20\% pre-flood to around 18\%
post-flood, indicating that freight partially diverted to road
hauliers and did not fully return to rail during the observation
window.

\begin{table}[H]
  \centering
  \caption{Descriptive statistics for the NATCOR corridor
    (120 monthly observations, 2015-01 to 2024-12).}
  \label{tab:natcor_stats}
  \small
  \renewcommand{\arraystretch}{1.2}
  \begin{tabular}{lrrrrr}
    \toprule
    \textbf{Variable} & \textbf{Mean} & \textbf{Std} &
      \textbf{Min} & \textbf{Median} & \textbf{Max} \\
    \midrule
    Rail\_Vol\_mt       & 3.87  & 0.42  & 2.81  & 3.94  & 4.73  \\
    Market\_Share\_Pct  & 19.4  & 2.8   & 13.1  & 19.8  & 24.6  \\
    Loco\_Avail\_Rate   & 0.641 & 0.027 & 0.573 & 0.643 & 0.695 \\
    Load\_Shedding\_Hrs & 48.3  & 64.1  & 0.0   & 17.4  & 214.6 \\
    Fuel\_Price\_Diesel & 15.82 & 3.41  & 10.50 & 15.60 & 23.80 \\
    Regional\_GDP\_Idx  & 108.2 & 7.9   & 91.3  & 108.6 & 121.4 \\
    Security\_Incidents & 63.1  & 6.2   & 54.0  & 61.5  & 80.0  \\
    Rainfall\_mm        & 116.8 & 33.7  & 57.1  & 117.6 & 179.3 \\
    \bottomrule
  \end{tabular}
\end{table}

\subsubsection{Temporal Exploratory Data Analysis}
\label{sec:eda}

Figure~\ref{fig:eda_all} shows the monthly Rail\_Vol\_mt series for
all six corridors over the full observation window.
Several features are immediately apparent.
First, NATCOR exhibits a clear downward level shift beginning in
April 2022, visible as a sustained deviation from the pre-flood
trend that persists through the end of the observation window.
Second, the other five corridors do not exhibit a comparable shift
at the same date, which is the fundamental identifying variation
that the SCM framework exploits.
Third, all corridors display a brief common dip in 2020--2021
corresponding to the COVID-19 disruption --- a common time shock
that SDID's time weights are specifically designed to downweight.

The NATCOR corridor exhibits mild positive seasonality in the
second and third quarters of each year, reflecting the agricultural
and container export peak.
Autocorrelation analysis confirms that the series is persistent
(autocorrelation at lag~1 $\approx 0.83$) and that shocks decay
slowly, making the pre-period balance requirement of SCM the
appropriate identification strategy.

\subsection{Data Preprocessing and Temporal Feature Engineering}
\label{sec:preprocessing}

The raw Excel workbook required several cleaning steps before
analysis, implemented in \texttt{data\_cleaning.ipynb} and
documented below.

\textbf{Step 1 --- Corridor name standardisation.}
Leading and trailing whitespace was stripped from the
\texttt{Corridor} field.
Six rows labelled \texttt{`NATCOR '} (with a trailing space) were
thereby correctly merged with the \texttt{`NATCOR'} group,
recovering missing NATCOR observations for January--March 2023.

\textbf{Step 2 --- Duplicate removal.}
Exact duplicate rows (identical Corridor and Date) were removed,
retaining the first occurrence.

\textbf{Step 3 --- Physiologically implausible values.}
\texttt{Loco\_Avail\_Rate} values exceeding 1.0 are impossible
(a rate is bounded above by 1) and were set to \texttt{NaN} for
subsequent interpolation.
\texttt{Avg\_Delay\_Hrs} values exceeding 200 hours were treated
as data entry errors and similarly nullified.
One \texttt{Rail\_Vol\_mt} entry of exactly zero (CAPE, November 2018)
was also nullified.

\textbf{Step 4 --- Date format repair.}
Six IRON\_ORE rows had dates in \texttt{YYYY-MM} format rather than
\texttt{YYYY-MM-DD}.
These were reconstructed from the \texttt{Year} and \texttt{Month}
columns with a day value of 01.

\textbf{Step 5 --- Linear interpolation.}
All remaining \texttt{NaN} values were filled by linear interpolation
within each corridor's time series, preserving chronological ordering.
Interpolation was applied to \texttt{Rail\_Vol\_mt},
\texttt{Road\_Vol\_Est\_mt}, \texttt{Avg\_Delay\_Hrs},
\texttt{Loco\_Avail\_Rate}, \texttt{Security\_Incidents}, and
\texttt{Rainfall\_mm}.
Static variables (\texttt{Fuel\_Price\_Diesel},
\texttt{Regional\_GDP\_Idx}, \texttt{Load\_Shedding\_Hrs}) had no
missing values.

\textbf{Step 6 --- Train/validation/test split.}
No random shuffling is applied at any stage.
The temporal ordering is always preserved.
For cross-validation within the SCM estimators, the pre-period
(87 months) is split chronologically: the first 60\% (52 months,
2015-01 to 2019-04) serves as the training fold and the remaining
40\% (35 months, 2019-05 to 2022-03) as the validation fold.
The post-period (33 months, 2022-04 to 2024-12) is held out
entirely and used only for treatment effect estimation.

\subsection{Ground-Truth Temporal Graph Construction}
\label{sec:groundtruth}

To evaluate the plausibility of the causal estimates and to provide
a structured prior for the estimand identification step, a
domain-informed ground-truth temporal graph is constructed.
Unlike the standard SCM setup --- which treats the counterfactual
as a weighted average of donors --- the temporal graph makes explicit
which variables are expected to causally influence NATCOR freight
volume, and over what lag.

The graph is represented as a set of directed edges of the form
$X_{t-k} \to Y_t$, where $X$ is a causal variable, $Y$ is
Rail\_Vol\_mt, and $k \geq 0$ is the lag in months.
The following edges are supported by domain knowledge and the
literature:

\begin{itemize}
  \item $\text{Loco\_Avail\_Rate}_{t} \to \text{Rail\_Vol}_{t}$
        (contemporaneous): locomotive availability directly
        determines operational capacity in the same month.
  \item $\text{Load\_Shedding\_Hrs}_{t-1} \to \text{Rail\_Vol}_{t}$
        (one-month lag): load-shedding disrupts rail terminal
        operations with a lag, as rolling stock scheduling adjustments
        take time to propagate.
  \item $\text{Security\_Incidents}_{t-1} \to \text{Rail\_Vol}_{t}$:
        cable theft and vandalism in month $t-1$ reduce available
        track capacity in month $t$.
  \item $\text{Regional\_GDP\_Idx}_{t} \to \text{Rail\_Vol}_{t}$:
        demand for freight services tracks regional economic activity.
  \item $\text{Flood\_Shock}_{t=\text{2022-04}} \to
        \text{Rail\_Vol}_{t}$ for $t \geq \text{2022-04}$:
        the treatment --- a discrete, externally-imposed shock
        that damaged NATCOR infrastructure.
  \item $\text{Rail\_Vol}_{t-1} \to \text{Rail\_Vol}_{t}$:
        strong autocorrelation in the freight series reflects
        contract-based shipping arrangements and shipper inertia.
\end{itemize}

Contemporaneous and lagged edges are kept separate in the graph.
This separation is important because the SCM framework implicitly
conditions on the donor pool to control for contemporaneous
confounders (variables that affect all corridors simultaneously,
such as load-shedding and GDP), while lagged effects --- in particular
the treatment effect itself --- are the primary object of estimation.

\subsection{Five SCM-Family Causal Estimators}
\label{sec:estimators}

Five estimators from the SCM family are applied to the corridor
panel.
All five share the same data, the same pre-period (87 months),
the same donor pool (five corridors), and the same treatment date
(April 2022).
They differ in how they construct the counterfactual.

\subsubsection{Base SCM}
\label{sec:base}

The standard constrained SCM solves Equation~\ref{eq:scm} via
Sequential Least Squares Programming (SLSQP), imposing
$w_j \geq 0$ and $\sum_j w_j = 1$.
No regularisation is applied.
The pre-period RMSPE serves as a direct measure of counterfactual
quality.
With five donors, the optimal weights are:
SOUTHCOR (0.463), NE\_CORRIDOR (0.241), CAPE (0.221),
NORTHCOR (0.076), and IRON\_ORE (0.000).
The pre-period RMSPE is 0.0261\,MT, confirming that the synthetic
NATCOR tracks the actual NATCOR very closely before the flood.

\subsubsection{Augmented SCM (ASCM)}
\label{sec:ascm_method}

ASCM adds a ridge bias-correction on top of the Base SCM
counterfactual (Section~\ref{sec:ascm}).
The ridge features include the five donor series, normalised time,
time-squared, and donor--time interaction terms.
The regularisation parameter $\alpha$ is selected by blocked
cross-validation over a 13-point grid spanning
$[10^{-5},\, 100]$.
The CV-selected $\alpha = 100$ indicates that a strong ridge penalty
is optimal --- the bias-correction term is heavily shrunk toward zero,
consistent with the already excellent pre-period fit of the Base SCM.

\subsubsection{Elastic Net SCM (EN-SCM)}
\label{sec:enscm_method}

EN-SCM replaces the hard sum-to-one constraint with an elastic net
penalty (Equation~\ref{eq:enscm}).
A two-dimensional blocked grid search over
$\alpha \in \{10^{-4}, 5\times10^{-4}, \ldots, 1.0\}$
(9 values) and
$\rho \in \{0.0, 0.1, 0.25, 0.5, 0.75, 0.9, 1.0\}$
(7 values) yields 63 candidate combinations.
The CV-selected parameters are $\alpha = 5\times10^{-4}$ and
$\rho = 0.25$, indicating a mild predominantly ridge-type penalty.
All five donors receive non-zero weights (active donors = 5),
with a Herfindahl--Hirschman Index (HHI) of 0.315 --- substantially
less concentrated than the Base SCM (HHI = 0.327).
The weight sum of 0.502 (less than 1) reflects the relaxed
constraint, which allows the model to discount donors that are
systematically larger or smaller than NATCOR.

\subsubsection{Matrix Completion SCM (MC-SCM)}
\label{sec:mcscm_method}

MC-SCM treats the $120 \times 6$ corridor-time panel as a
low-rank matrix and imputes the 33 missing post-treatment NATCOR
entries via Singular Value Thresholding
(Equation~\ref{eq:mcscm}).
The nuclear norm penalty $\lambda$ is selected by holding out the
last 25 pre-period NATCOR observations as a masked validation set
and searching over 14 candidate values.
The CV-selected $\lambda = 0.2$ yields an effective matrix rank of 4,
meaning that four latent factors explain essentially all of the
cross-corridor covariance structure.
The pre-period RMSPE of MC-SCM is exactly 0 by construction (all
pre-period entries are observed), and the post-period RMSPE of
0.167\,MT is the lowest among all five estimators.

\subsubsection{Synthetic Difference-in-Differences (SDID)}
\label{sec:sdid_method}

SDID solves for unit weights $\boldsymbol{\lambda}$ (identical to
the Base SCM weights) and time weights $\boldsymbol{\omega}$
separately (Equation~\ref{eq:sdid}).
The time weight regularisation parameter $\lambda_\omega$ is
selected from a sensitivity grid; the near-zero optimal value
($\lambda_\omega = 10^{-6}$) indicates that the time weights are
nearly uniform across pre-period months, reflecting a relatively
stable pre-period donor structure.
The effective number of pre-period months contributing to the
estimate is 60.01 out of 87 --- high entropy weights that do not
concentrate on a small subset of pre-flood months.
Standard errors are computed by the jackknife, dropping one
pre-treatment year at a time (eight jackknife folds, 2015--2022).

\subsection{Discovery Evaluation Metrics}
\label{sec:metrics}

Each estimator is evaluated on the following metrics:

\textbf{Pre-period RMSPE.}
The root mean squared prediction error between the synthetic control
and the actual NATCOR series in the pre-period:
\begin{equation}
  \text{RMSPE}_{\text{pre}} = \sqrt{\frac{1}{T_0}
    \sum_{t=1}^{T_0}\!\left(Y_{1t} - \hat{Y}_{1t}\right)^2}
\end{equation}
A lower pre-period RMSPE indicates a tighter fit and a more credible
counterfactual.

\textbf{Post-period RMSPE.}
The analogous metric computed over the post-treatment period.
A high post-period RMSPE, relative to the pre-period RMSPE, is
evidence that the treatment caused a divergence between the actual
outcome and its counterfactual.

\textbf{Post/Pre RMSPE ratio.}
The ratio of post-period to pre-period RMSPE.
A ratio substantially greater than 1 indicates that the deviation
in the post-period is larger than typical pre-period noise,
consistent with a genuine treatment effect.
This ratio is also the basis for the placebo-based inference described
below.

\textbf{Cumulative treatment effect.}
$\hat{\tau}_{\text{cum}} = \sum_{t=T_0+1}^{T} \hat{\tau}_t$, the
sum of monthly gap estimates over the 33 post-period months.
This quantity represents the total freight volume loss (or gain) in
millions of tonnes attributable to the flood shock.

\textbf{Average monthly treatment effect.}
$\bar{\tau} = \hat{\tau}_{\text{cum}} / (T - T_0)$, the average
causal effect per month.

\textbf{Post-period MAPE.}
The mean absolute percentage error over the post-period, which
provides a scale-free measure of counterfactual accuracy.

\subsection{Robustness and Stability Assessment}
\label{sec:robustness}

\subsubsection{In-Time Placebo Tests}
\label{sec:placebo}

For each estimator, five fake treatment dates are assigned within
the pre-period: 2017-01, 2018-06, 2019-07, 2020-04, and 2021-01.
For each placebo date $t^*$, the data prior to $t^*$ is used as the
``pre-period'' and the data from $t^*$ to 2022-03 serves as the
``post-period'', with the actual post-treatment data excluded.
The estimator is re-fitted and a placebo RMSPE ratio is computed.

If the true post/pre ratio (using the actual April 2022 treatment
date) exceeds all five placebo ratios, this is strong evidence that
the detected divergence is specific to the flood event and not an
artefact of the modelling approach.
If some placebo ratios exceed the true ratio, the inference is
weakened proportionally.

\subsubsection{Leave-One-Out Donor Stability}
\label{sec:loo}

Each of the five donor corridors is removed from the donor pool in
turn, and the estimator is re-fitted on the remaining four donors.
The spread of cumulative effect estimates across the five
leave-one-out runs quantifies how sensitive the estimate is to
any single donor corridor.
A spread below 0.5\,MT is considered stable; a spread above 2\,MT
indicates fragility.

\subsubsection{Hyperparameter Sensitivity}
\label{sec:sensitivity}

For ASCM, the full $\alpha$ grid ($10^{-5}$ to $10^2$) is explored
and the resulting cumulative effect and direction are reported.
For EN-SCM, the direction consistency across all 63 grid combinations
is reported (i.e., what fraction of hyperparameter settings yield a
negative estimated effect).
For MC-SCM, the $\lambda$ sensitivity table covers 14 values.
For SDID, the jackknife standard error and 95\% confidence interval
are reported.

\subsection{Causal Estimand Identification}
\label{sec:estimand}

The causal estimand of interest is the \textbf{Average Treatment
Effect on the Treated} (ATT):
\begin{equation}
  \text{ATT} = \mathbb{E}\!\left[Y_{1t}(1) - Y_{1t}(0)
    \;\middle|\; \text{NATCOR treated}\right]
    \quad \text{for } t > T_0
\end{equation}
where $Y_{1t}(1)$ is the potential outcome under treatment (the
observed NATCOR volume post-April 2022) and $Y_{1t}(0)$ is the
potential outcome in the absence of treatment (the counterfactual
constructed by the SCM).

The treatment variable is the April 2022 KZN flood shock, modelled
as a binary indicator $D_t = \mathbf{1}[t \geq \text{2022-04}]$
for the NATCOR corridor.
The outcome variable is monthly Rail\_Vol\_mt.
The adjustment set consists of the five donor corridors, which are
assumed to share common time-varying confounders (load-shedding,
fuel prices, GDP, security conditions) with NATCOR but to be
unaffected by the KZN flood shock.

The temporal back-door logic is as follows.
Let $\mathbf{C}_t$ denote the vector of common confounders at time
$t$ (load-shedding, GDP, etc.).
These confounders affect both NATCOR and the donors
contemporaneously.
By constructing a synthetic control whose pre-period trajectory
matches NATCOR --- and which continues to be driven by the same
confounders after April 2022 --- the SCM implicitly conditions on
$\mathbf{C}_t$ in the post-period.
The identification assumption is that, conditional on this
pre-period balance, the post-period divergence between NATCOR and
its synthetic control is attributable to the flood shock and to no
other corridor-specific event.

Potential threats to this assumption include:
\begin{itemize}
  \item Corridor-specific deterioration in locomotive availability
        that is unrelated to the flood but coincides with it.
  \item Changes in Transnet policy (e.g., prioritisation of the
        Iron Ore line) that affect NATCOR but not donors after 2022.
  \item Continued security incident escalation that is more severe
        on NATCOR than donors in the post-period.
\end{itemize}
These threats are partially mitigated by the multi-estimator design
and the robustness tests, but cannot be fully eliminated with
observational data alone.

\subsection{Causal Effect Estimation}
\label{sec:estimation}

The primary causal effect estimate is the monthly treatment gap
series $\hat{\tau}_t = Y_{1t} - \hat{Y}_{1t}^{(m)}$, where
$\hat{Y}_{1t}^{(m)}$ is the counterfactual produced by estimator
$m \in \{\text{Base SCM, ASCM, EN-SCM, MC-SCM, SDID}\}$.

For SDID, the ATT has a closed-form expression
(Equation~\ref{eq:sdid}) and is accompanied by a jackknife
standard error:
\begin{equation}
  \widehat{\text{SE}}_{\text{JK}} = \sqrt{\frac{N-1}{N}
    \sum_{k=1}^{N}\!\left(\hat{\tau}^{(-k)} - \bar{\tau}\right)^2}
\end{equation}
where $\hat{\tau}^{(-k)}$ is the ATT estimate when the $k$-th
pre-treatment year is removed and $\bar{\tau}$ is the full-sample
estimate.
A 95\% confidence interval is computed as
$\bar{\tau} \pm 1.96 \cdot \widehat{\text{SE}}_{\text{JK}}$.

For MC-SCM, a block bootstrap with 200 resamples and a block size
of 6 months is used to construct a 95\% confidence interval for
the cumulative effect.

For Base SCM, ASCM, and EN-SCM, inference is based on the
placebo-ratio comparison described in Section~\ref{sec:placebo},
which is the standard Fisher randomisation test for SCM.

The main causal estimator for the purposes of policy interpretation
is the \textbf{consensus estimate}: the arithmetic mean of the five
cumulative effect estimates, which is robust to the specific
modelling assumptions of any single method.

\subsection{Refutation, Sensitivity, and Reproducibility}
\label{sec:refutation}

\subsubsection{Refutation Tests}
\label{sec:refutation_tests}

Three formal refutation tests are conducted.

\textbf{Placebo treatment date.}
The treatment date is shifted to April 2020 (the COVID-19
disruption) and April 2021 (a placebo with no known major shock).
If the estimators detect large effects at these placebo dates, it
suggests that they are responding to systematic trends rather than
the flood.

\textbf{Subset-by-period refuter.}
The analysis is repeated using only the pre-COVID pre-period
(2015-01 to 2019-12) and re-estimated using a 2022-04 treatment
date.
Stable estimates across this constrained pre-period confirm that the
COVID disruption period does not drive the results.

\textbf{Alternative lag window.}
The treatment date is shifted by $\pm$3 months (to 2022-01 and
2022-07) to assess whether the estimated effect is sensitive to the
exact assumed start of the treatment.

\subsubsection{Reproducibility}
\label{sec:repro}

All analyses are implemented in Python~3.10.
Key libraries and versions: \texttt{pandas}~1.5, \texttt{numpy}~1.24,
\texttt{scipy}~1.10, \texttt{scikit-learn}~1.2,
\texttt{matplotlib}~3.7.
All random seeds are fixed at 42 where applicable.
No random shuffling is applied to any time-series split.
The cleaned dataset, all notebooks, and the figure output directory
are archived in the thesis repository, enabling full replication of
every result reported in Section~\ref{sec:results}.
